% chktex-file 3
% chktex-file 8
% chktex-file 35
\section{Boundary thermodynamics and power scaling}\label{sec:thermodynamics-power}

The completed dark ledger supplies the positive modular entropy debt of the
operating branch,
\begin{equation}
  \Delta_{\mathrm{mod}}
  =\frac{c_{\mathrm{dark}}^{\mathrm{comp}}}{24}
  =\frac{1197103}{8704080}
  =\SimOutputEntropyDebt.
  \label{eq:thermodynamic-entropy-debt}
\end{equation}
Within the kinematic ansatz of Sec.~\ref{sec:shift-derivation}, the associated
shift scale is
\begin{equation}
  v_{\mathrm{eff}}
  =c\exp\left(\frac{\Delta_{\mathrm{mod}}}{2}\right),
  \qquad
  \frac{v_{\mathrm{eff}}}{c}
  =\SimOutputWarpVelocity.
  \label{eq:thermodynamic-effective-velocity}
\end{equation}
This exponential is a constitutive relation of the extension, not a
consequence of the Einstein equations alone.

For a driven boundary register, a minimal relaxation model separates the
control injection rate \(\mathcal F\) from passive dark-sector dissipation,
\begin{equation}
  \frac{d\Delta_{\mathrm{mod}}}{dt}
  =\mathcal F(P_{\mathrm{op}},\theta,R_{\mathrm{local}})
  -\frac{\Gamma_{\mathrm{lock}}}{24}\Delta_{\mathrm{mod}}.
  \label{eq:entropy-debt-rate-equation}
\end{equation}
For constant drive, the solution is
\begin{equation}
  \Delta_{\mathrm{mod}}(t)
  =\Delta_{\mathrm{mod}}^{\mathrm{ss}}
  +\left[
    \Delta_{\mathrm{mod}}(0)-\Delta_{\mathrm{mod}}^{\mathrm{ss}}
  \right]
  \exp\left(-\frac{\Gamma_{\mathrm{lock}}t}{24}\right),
  \qquad
  \Delta_{\mathrm{mod}}^{\mathrm{ss}}
  =\frac{24\mathcal F}{\Gamma_{\mathrm{lock}}}.
  \label{eq:entropy-debt-relaxation-solution}
\end{equation}
The benchmark fixes the steady state to
Eq.~\eqref{eq:thermodynamic-entropy-debt}; a microscopic derivation of
\(\Gamma_{\mathrm{lock}}\) remains outside the present finite-register
simulation.

\subsection{Calibrated operational power law}

The operational model uses the Planck-power scale multiplied by the modular
coupling and a geometric suppression factor,
\begin{equation}
  P_{\mathrm{op}}(R_{\mathrm{local}})
  =\frac{c^5}{G_N}
  \frac{\Delta_{\mathrm{mod}}}{24\pi}
  \left(\frac{r_g}{R_{\mathrm{local}}}\right)^2.
  \label{eq:operational-power-fundamental}
\end{equation}
The length \(r_g\) is not fixed by the dimensionless dark ledger.  Therefore
the quoted power benchmark constitutes an independent calibration.  Solving
Eq.~\eqref{eq:operational-power-fundamental} for that scale gives
\begin{equation}
  r_g^{\mathrm{cal}}
  =R_0
  \left[
    \frac{24\pi G_N P_0}
    {c^5\Delta_{\mathrm{mod}}}
  \right]^{1/2}
  =\SimOutputPowerScaleRadius\ \mathrm{m},
  \label{eq:calibrated-power-length}
\end{equation}
where
\begin{equation}
  R_0=10\ \mathrm{m},
  \qquad
  P_0=\SimOutputOperationalPowerMW\ \mathrm{MW}.
  \label{eq:power-benchmark-pair}
\end{equation}
Substitution produces the inverse-area law used by the executable engine,
\begin{equation}
  P_{\mathrm{op}}(R_{\mathrm{local}})
  =\SimOutputOperationalPowerMW
  \left(\frac{10\ \mathrm{m}}{R_{\mathrm{local}}}\right)^2
  \mathrm{MW}.
  \label{eq:operational-power-scaling}
\end{equation}
In particular,
\begin{equation}
  P_{\mathrm{op}}(10\ \mathrm{m})
  =\SimOutputPowerReq\ \mathrm{MW}.
  \label{eq:ten-meter-power-result}
\end{equation}
Equations~\eqref{eq:operational-power-fundamental}--\eqref{eq:ten-meter-power-result}
make the normalization transparent: the
\(142.08\,\mathrm{MW}\) value is conditional on
\(r_g=r_g^{\mathrm{cal}}\), whereas the \(R^{-2}\) dependence is the actual
prediction of the stated power ansatz.